\documentclass{article}
\usepackage{amsmath, amssymb , amsthm, graphicx,mathtools}

\begin{document}

\section*{Question 1}

~

\begin{proof}
    ~
    \begin{enumerate}
        \item \begin{align*}
            &\binom{n-1}{k}+\binom{n-1}{k-1}\\
            =&\frac{(n-1)!}{k!(n-1-k)!}+\dfrac{(n-1)!}{(k-1)!(n-k)!}\\
            =&\dfrac{(n-k)(n-1)!}{k!(n-k)!}+\dfrac{k(n-1)!}{k!(n-k)!}\\
            =&\dfrac{(n-k)(n-1)!+k(n-1)!}{k!(n-k)!}\\
            =&\dfrac{n!}{k!(n-k)!}\\
            =&\binom{n}{k}
        \end{align*}
        \item \begin{enumerate}
            \item $\binom{n}{k}=\binom{n-1}{k}+\binom{n-1}{k-1}=\binom{n-2}{k}+\binom{n-2}{k-1}+\binom{n-1}{k-1}=...$
            \item by this iteration on every $\binom{i}{k}$: $\binom{n}{k}=\binom{k}{k}+\sum_{j=k}^{n-1}\binom{j}{k-1}$
            \item $\binom{k}{k}=1=\binom{k-1}{k-1}$
            \item Substitution: $\binom{n}{k}=\binom{k-1}{k-1}+\sum_{j=k}^{n-1}\binom{j}{k-1}=\sum_{j=k-1}^{n-1}\binom{j}{k-1}$
        \end{enumerate}
    \end{enumerate}
\end{proof}

\newpage

\section*{Question 2}

~

\begin{proof}
    ~
    \begin{enumerate}
        \item Set $AD\parallel A'D''$
        \item $\angle OAB=\angle OA'B'',\angle ABO=\angle A'B''O, \angle OBC=\angle OB''C'', \angle OCD=\angle OC''D'', \angle OCB= \angle OC''B',\angle ODA= \angle OD''A'$
        \item $\triangle ABO\sim \triangle A'B''O,\triangle CDO\sim \triangle C''D''O,\triangle BCO\sim \triangle B''C''O,\triangle ADO\sim \triangle A'D''O$
        \item $\dfrac{AB}{A'B''}=\dfrac{AO}{A'O}=\dfrac{AD}{A'D''}$, $\dfrac{BC}{B''C''}=\dfrac{OD}{OD''}=\dfrac{CD}{C''D''}$
        \item $\dfrac{(A,B;C,D)}{(A',B'';C'',D'')}=\dfrac{AB\cdot CD\cdot A'D''\cdot C''B''}{AD\cdot CB\cdot A'B''\cdot C''D''}$
        \item $\dfrac{AB\cdot CD\cdot A'D''\cdot C''B''}{AD\cdot CB\cdot A'B''\cdot C''D''}=\dfrac{AB}{A'B''}\cdot\dfrac{CD}{C''D''}\cdot\dfrac{A'D''}{AD}\cdot\dfrac{C''B''}{CB}=\dfrac{AB}{A'B''}\cdot\dfrac{CD}{C''D''}\cdot\dfrac{1}{\frac{AB}{A'B''}}\cdot\dfrac{1}{\frac{CD}{C''D''}}=1$
        \item So $(A,B;C,D)=(A',B'';C'',D'')$
        \item $\dfrac{(A',B';C',D')}{(A',B'';C'',D'')}=\dfrac{A'B'\cdot C'D'\cdot A'D''\cdot C''B''}{A'D'\cdot C'B'\cdot A'B''\cdot C''D''}$
        \item \begin{enumerate}
            \item Menelaus theorem on $\triangle AC'C''$ with $OD'$: $\dfrac{A'D'}{D'C'}\cdot\dfrac{C'O}{C''O}\cdot\dfrac{C''D''}{A'D''}=1$
            \item Menelaus theorem on $\triangle AC'C''$ with $OB'$: $\dfrac{A'B'}{B'C'}\cdot\dfrac{C'O}{C''O}\cdot\dfrac{C''B''}{A'B''}=1$
        \end{enumerate}
        \item $\dfrac{D'C'}{A'D'}\cdot\dfrac{C''O}{C'O}\cdot\dfrac{A'D''}{C''D''}\cdot \dfrac{A'B'}{B'C'}\cdot\dfrac{C'O}{C''O}\cdot\dfrac{C''B''}{A'B''}=1$
        \item $\dfrac{A'B'\cdot C'D'\cdot A'D''\cdot C''B''}{A'D'\cdot C'B'\cdot A'B''\cdot C''D''}=1$
        \item So $(A',B';C',D')=(A',B'';C'',D'')$
        \item $(A,B;C,D)=(A',B';C',D')$
    \end{enumerate}
\end{proof}

\newpage

\section*{Question 3}

~

\begin{proof}
    ~
    \begin{enumerate}
        \item Set the radius of the large sphere $r$ and the random distance above the equator of the sphere $-\dfrac{h}{2}\leq y\leq \dfrac{h}{2}$
        \item Radius of the cylinder is $\sqrt{r^2-\left(\dfrac{h}{2}\right)^2}$ and radius for cross-section at $y$ is $\sqrt{r^2-y^2}$
        \item The difference in area of these two is the area of the ring at height $y$: $\pi\left(-r^2+\left(\dfrac{h}{2}\right)^2+r^2-y^2\right)=\pi\left(\left(\dfrac{h}{2}\right)^2-y^2\right)$
        \item Consider a sphere with radius $\dfrac{h}{2}$, the radius of one circle of the sphere at height $y$ is $\sqrt{\left(\dfrac{h}{2}\right)^2-y^2}$
        \item The area of the circle is $\pi\left(\left(\dfrac{h}{2}\right)^2-y^2\right)$
        \item The two areas are the same.
        \item By Cavalieri’s principle, the two volumes with equal-sized of corresponding cross-sections are equal-volumed.
        \item The ring and the $\dfrac{h}{2}$-radius sphere have the same volume.
        \item The volume of the ring is independent of the radius $r$.
    \end{enumerate}
\end{proof}

\newpage

~

\section*{Question 4}

~

\begin{enumerate}
    \item $P(a,b)$ is the probability Cain wins when he needs $a$ points and Abel needs $b$ points.
    \item $P(a,b)= 0.5P(a-1,b)+0.5P(a,b-1)$, $P(0,b)=1$, $P(a,0)=0$
    \item \begin{align*}
        &P(3,4)\\
        =&0.5P(2,4)+0.5P(3,3)\\
        =&0.25P(1,4)+0.5P(2,3)+0.25P(3,2)\\
        =&0.125P(0,4)+0.375P(1,3)+0.375P(2,2)+0.125P(3,1)\\
        =&0.125+0.1875P(0,3)+0.375P(1,2)+0.25P(2,1)+0.0625P(3,0)\\
        =&0.125+0.1875+0.1875P(0,2)+0.3125P(1,1)+0.125P(2,0)\\
        =&0.125+0.1875+0.1875+0.15625P(0,1)+0.15625P(1,0)\\
        =&0.125+0.1875+0.1875+0.15625\\
        =&0.65625\\
    \end{align*}
    \item $0.65625\times 64=42$, $64-42=22$
    \item 42 apples for Cain and 22 apples for Abel.
\end{enumerate}

\newpage

\section*{Question 5}

~

\begin{proof}
    ~
    \begin{enumerate}
        \item The proof is same as: $\forall n>2\in\mathbb{Z}: \exists p_\text{odd}, p_\text{odd}|n\lor 4|n$
        \item For all integers greater than 2, it can only be the case that some odd primes divide it or no odd primes divide it.
        \item If some odd primes divide it, then $\exists p_\text{odd}:p_\text{odd}|n$
        \item If no odd primes divide it, then $n$ can only be the form of $2^k$ for some integer $k$.
        \item Since $n>2$, $k\geq2$
        \item $4=2^2$ must divides $2^k,k\geq 2$, so $4|n$
        \item Combining the two cases: $\forall n>2\in\mathbb{Z}: \exists p_\text{odd}, p_\text{odd}|n\lor 4|n$
    \end{enumerate}
\end{proof}

\newpage

\section*{Question 6}

~

\begin{proof}
    ~
    \begin{enumerate}
        \item Base Case ($n=1$): $y'=\sum_{k=0}^1 \binom{1}{k} u^{(k)}v^{(1-k)}=u'v+uv'$
        \item Induction Hypothesis ($n=m$): $y^{(m)}=\sum_{k=0}^m \binom{m}{k} u^{(k)}v^{(m-k)}$
        \item Induction Step ($n=m+1$):
        \begin{enumerate}
            \item $y^{(m+1)}=\frac{d}{dx}y^{(m)}=\frac{d}{dx}\sum_{k=0}^m \binom{m}{k} u^{(k)}v^{(m-k)}$ 
            \item $\frac{d}{dx}\sum_{k=0}^m \binom{m}{k} u^{(k)}v^{(m-k)}=\sum_{k=0}^m \binom{m}{k} \left(u^{(k+1)}v^{(m-k)}+u^{(k)}v^{(m-k+1)}\right)$
            \item $\sum_{k=0}^m \binom{m}{k} \left(u^{(k+1)}v^{(m-k)}+u^{(k)}v^{(m-k+1)}\right)=u^{m+1}v+\sum_{k=1}^m \binom{m}{k-1}u^{(k)}v^{(m-k+1)}+ \sum_{k=1}^m \binom{m}{k}u^{(k)}v^{(m-k+1)}+ uv^{(m+1)}$
            \item $u^{m+1}v+\sum_{k=1}^m \binom{m}{k-1}u^{(k)}v^{(m-k+1)}+ \sum_{k=1}^m \binom{m}{k}u^{(k)}v^{(m-k+1)}+ uv^{(m+1)}= u^{m+1}v+\sum_{k=1}^m \left(\binom{m}{k-1}+\binom{m}{k}\right)u^{(k)}v^{(m-k+1)}+ uv^{(m+1)}$
            \item $\binom{m}{k-1}+\binom{m}{k}=\binom{m+1}{k}$
            \item $u^{m+1}v+\sum_{k=1}^m \left(\binom{m}{k-1}+\binom{m}{k}\right)u^{(k)}v^{(m-k+1)}+ uv^{(m+1)}=\binom{m+1}{m+1}u^{m+1}v+\sum_{k=1}^m \binom{m+1}{k}u^{(k)}v^{(m-k+1)}+ \binom{m+1}{0}uv^{(m+1)}= \sum_{k=0}^m \binom{m+1}{k}u^{(k)}v^{(m-k+1)}$
            \item So $y^{(m+1)}= \sum_{k=0}^m \binom{m+1}{k}u^{(k)}v^{(m-k+1)}$
        \end{enumerate}
        \item By induction, $\forall n\in\mathbb{N}, y^{(n)}= \sum_{k=0}^m \binom{n}{k}u^{(k)}v^{(n-k)}$
    \end{enumerate}
\end{proof}

\newpage

\section*{Question 7}

~

\begin{enumerate}
    \item Fermat’s Method of Adequality
    \begin{enumerate}
        \item $x\coloneqq 3+m, y\coloneqq 4+n$
        \item $(3+m)^2+(4+n)^2=25$
        \item $6m+8n+m^2+n^2=25$, where $m^2+n^2$ are negligible
        \item $6m+8n\approx 0$
        \item $\dfrac{n}{m}=-\dfrac{3}{4}$
        \item Tangent is $-\dfrac{3}{4}$
        \item $l: y=-\dfrac{3}{4}x+\dfrac{25}{4}$
    \end{enumerate}
    \item Descartes’ Strategy
    \begin{enumerate}
        \item Slope of radius is $\frac{4}{3}$
        \item Normal line is $y=-\dfrac{4}{3}x$
        \item Tangent slope perpendicular to normal slope is $-\dfrac{3}{4}$
        \item Tangent line $y=-\dfrac{3}{4}x+\dfrac{25}{4}$
    \end{enumerate}
    \item Newton’s Fluxions
    \begin{enumerate}
        \item Differentiation: $2x\dot{x}+2y\dot{y}=0$
        \item $\dfrac{\dot{y}}{\dot{x}}=-\dfrac{x}{y}$
        \item $x=3,y=4$
        \item $\dfrac{\dot{y}}{\dot{x}}=-\dfrac{3}{4}$
        \item $l: y=-\dfrac{3}{4}x+\dfrac{25}{4}$
    \end{enumerate}
\end{enumerate}
\end{document}